
\noindent\textbf{命题~\ref{prop:7-1-27} 的证明。}
由假设，存在有界线性映射
\[
Q_0\colon
\frac{\mathcal E^{0,p}_{((\mathfrak t_1,\ell_1),w_1,\vec p^{\,1})}}
{\mathcal O_1}
\oplus
\frac{\mathcal E^{0,p}_{((\mathfrak t_0,\ell_0),w_0,\vec p^{\,0})}}
{\mathcal O_0}
\longrightarrow\mathcal X_{\mathrm{small}},
\]
使得
\begin{equation}
\pi\circ D_{(w_1,w_0)}\bar\partial\circ Q_0
=\operatorname{id}.
\tag{7.1.30}\label{eq:7-1-30}
\end{equation}
这里，$D_{(w_1,w_0)}\bar\partial$
是
$D_{((\mathfrak t_1,\ell_1),w_1,\vec p^{\,1})}\bar\partial
\oplus
D_{((\mathfrak t_0,\ell_0),w_0,\vec p^{\,0})}\bar\partial$
在 $\mathcal X_{\mathrm{small}}$ 上的限制。

设 $V$ 属于胶合后纤维
$\mathcal E^{0,p}_{((\mathfrak t',\ell'_T),w,\vec p')}$，
以 $\bar V$ 表示它在该纤维模
$\mathcal O_1\oplus\mathcal O_0$ 后的等价类。
取 $2<S<5T-2$，定义近似右逆
\translatornote{原书式 (7.1.31) 把 $J_{*,0}^{-S}(\chi_0^+V)$ 放进 $Q_0$ 定义域的分支 1 槽，把 $J_{*,1}^{S}(\chi_0^-V)$ 放进分支 0 槽，且将 $V$ 与 $\bar V$ 误置于断开分支的直和空间。译文依前述 $J$ 的源靶和式 (7.1.30) 的型别先写分支 1、再写分支 0；后文同型元组亦如此。见疑误 A-0420。}
\begin{equation}
\begin{aligned}
&\widetilde Q_{((\mathfrak t',\ell'_T),w)}(\bar V)\\
&\quad=(I^T\circ Q_0)\bigl(
J_{*,1}^{S}(\chi_0^-V)
\oplus J_{*,0}^{-S}(\chi_0^+V)\bigr).
\end{aligned}
\tag{7.1.31}\label{eq:7-1-31}
\end{equation}
其中 $I^T$ 如式~\eqref{eq:7-1-23}。
由于 $\mathcal O_1,\mathcal O_0$ 中的截面支集
均不与新边 $L_{e(\mathrm{new})}$ 相交，
上述两项分别在
$\mathcal E^{0,p}_{((\mathfrak t_1,\ell_1),w_1,\vec p^{\,1})}/\mathcal O_1$
和
$\mathcal E^{0,p}_{((\mathfrak t_0,\ell_0),w_0,\vec p^{\,0})}/\mathcal O_0$
中的等价类只依赖 $\bar V$。
而且，只要 $2<S<5T-2$，
式~\eqref{eq:7-1-31} 的右端就不依赖 $S$：
$\chi_0^+V$ 在 $[-5T,-2]\times[0,1]$ 上为零，
$\chi_0^-V$ 在 $[2,5T]\times[0,1]$ 上为零。

研究 $\widetilde Q_{((\mathfrak t',\ell'_T),w)}$
需要下面两个关键估计引理。

\begin{lemma}{7.1.32}\label{lem:7-1-32}
设
$\mathbf V=((0,V_1,\vec v^{\,1}),(0,V_0,\vec v^{\,0}))
\in\mathcal X_{\mathrm{small}}$，则
\[
\begin{aligned}
&\bigl\|(D_w\bar\partial\circ I^T)(\mathbf V)
 -(J_{1,*}^{T}+J_{0,*}^{-T})
   D_{(w_1,w_0)}\bar\partial(\mathbf V)\bigr\|\\
&\qquad\leq
C\bigl(e^{-cT}
 +\operatorname{dist}(w,w'_0)
 +\operatorname{dist}(w,w'_1)\bigr)\|\mathbf V\|.
\end{aligned}
\]
\end{lemma}

\begin{proof}
$I^T$ 和 $J$ 与两个线性化算子不交换，
有两个原因。其一，须沿最短测地线
把 $w'_0,w'_1$ 处的截面平行移动到 $w$；
由此产生的误差由引理右端的两个距离项控制。
其二，也是更重要的原因，是这里用了截断函数和加权范数。
式~\eqref{eq:7-1-23} 中相关的两项是
\[
\chi_{-T}^+(\tau)
 \bigl(\operatorname{Pal}(V_0(\tau,t))-v\bigr),
\qquad
\chi_{+T}^-(\tau)
 \bigl(\operatorname{Pal}(V_1(\tau,t))-v\bigr).
\]
两项的处理相同，故只看第一项。
$\chi_{-T}^+$ 在
$[-T-1,-T+1]\times[0,1]$ 以外局部常值；
范数中需要估计的交换子项于是支撑在该条带上。
算子定义域中 $V_0$ 的权重在此处为
$Ce^{\delta(T+5T)}$。
这里将
$X(\mathfrak t_0,\ell_0)$ 的端部
识别为 $(-\infty,5T]\times[0,1]$，
而非 $(-\infty,0]\times[0,1]$，
故权重为 $e^{\delta|5T-\tau|'}$，
而非 $e^{\delta|\tau|'}$；
见定义~\ref{def:7-1-3}。
另一方面，值域的权重为 $Ce^{\delta(5T-T)}$。
因此两种权重之比给出所需的 $Ce^{-cT}$ 衰减；
见图~\ref{fig:7-1-6}。
\end{proof}

\par\medskip\noindent\begin{minipage}{\textwidth}
\centering
\setlength{\unitlength}{0.68mm}
\begin{picture}(145,80)
\linethickness{0.55pt}
\qbezier(43,73)(50,38)(69,24)
\qbezier(69,24)(78,15)(91,9)
\qbezier(39,10)(54,15)(65,28)
\qbezier(65,28)(75,38)(83,33)
\qbezier(83,33)(96,15)(127,9)
\put(50,48){\circle{15}}
\put(54,17){\circle{15}}
\put(17,59){\vector(0,-1){27}}
\put(2,43){$\scriptstyle\times e^{-cT}$}
\put(2,65){$\scriptstyle e^{\delta|5T-\tau|}$}
\put(5,15){$\scriptstyle e^{\delta|5T+\tau|}$}
\put(88,66){$\scriptstyle X(\mathfrak t_0,\ell_0)\text{ 上的权重}$}
\put(88,36){$\scriptstyle X(\mathfrak t',\ell'_T)\text{ 上的权重}$}
\end{picture}
\par\smallskip
{\renewcommand{\thefigure}{7.1.6}
\refstepcounter{figure}
{\normalfont\bfseries 图~\thefigure}\label{fig:7-1-6}}
\end{minipage}\par\medskip

\begin{lemma}{7.1.33}\label{lem:7-1-33}
存在与 $T$ 无关的常数 $C$，使
\[
\bigl\|\widetilde Q_{((\mathfrak t',\ell'_T),w)}
(\bar V)\bigr\|_{1,p,\alpha}
<C\|\bar V\|_{0,p,\alpha}.
\]
\end{lemma}

\begin{proof}
先证
\begin{equation}
\begin{aligned}
&\bigl\|I^T((0,V_1,\vec v^{\,1}),
                 (0,V_0,\vec v^{\,0}))\bigr\|_{1,p,\alpha}\\
&\qquad<C\bigl(
\|(V_1,\vec v^{\,1})\|_{1,p}
+\|(V_0,\vec v^{\,0})\|_{1,p}\bigr),
\end{aligned}
\tag{7.1.34}\label{eq:7-1-34}
\end{equation}
其中 $C$ 与 $T$ 无关。
只需估计左端在
$L_{e(\mathrm{new})}\cong[-5T,5T]\times[0,1]$
上的范数；其余部分的估计是直接的。
令
$V'=I^T((0,V_1,\vec v^{\,1}),
         (0,V_0,\vec v^{\,0}))$。
由式~\eqref{eq:7-1-23}，在 $|\tau|\leq1$ 上有
\[
V'(\tau,t)=\operatorname{Pal}\bigl(
\operatorname{Pal}(V_0(\tau,t))
+\operatorname{Pal}(V_1(\tau,t))-v\bigr).
\]
事实上，$-T+1<\tau<T-1$ 时也成立（这里 $T>2$）。
由引理~\ref{lem:7-1-5} 中的范数定义，
\[
\begin{aligned}
|V'(0,1/2)-v|<Ce^{-\delta c}
\bigl(&\|(0,V_1,\vec v^{\,1})\|_{1,p}\\
&+\|(0,V_0,\vec v^{\,0})\|_{1,p}\bigr).
\end{aligned}
\]
于是只需估计
\begin{equation}
\int_{L_{e(\mathrm{new})}}
\alpha\bigl(|\nabla V'|+|V'-v|\bigr)^p.
\tag{7.1.35}\label{eq:7-1-35}
\end{equation}
它由下面四个积分之和乘以常数 $C$ 控制：
\[
\begin{aligned}
&\int_{-5T}^{5T}\!d\tau\int_0^1
\alpha\bigl|\chi_{-T}^+(V_0-v)\bigr|^pdt
+\int_{-5T}^{5T}\!d\tau\int_0^1
\alpha\bigl|\chi_{+T}^-(V_1-v)\bigr|^pdt,\\
&\int_{-5T}^{5T}\!d\tau\int_0^1
\alpha\bigl|\nabla(\chi_{-T}^+(V_0-v))\bigr|^pdt
+\int_{-5T}^{5T}\!d\tau\int_0^1
\alpha\bigl|\nabla(\chi_{+T}^-(V_1-v))\bigr|^pdt.
\end{aligned}
\]
这里的 $\alpha$ 可由定义两个分支范数
$\|(V_0,\vec v^{\,0})\|_{1,p}$、
$\|(V_1,\vec v^{\,1})\|_{1,p}$ 所用的权重控制，
从而得到式~\eqref{eq:7-1-34}。

为完成本引理，还须证明
\begin{align}
\bigl\|J_{*,0}^{-T}(\chi_0^+V)\bigr\|_{0,p}
&\leq\|V\|_{0,p,\alpha},
\tag{7.1.36.1}\label{eq:7-1-36-1}\\
\bigl\|J_{*,1}^{T}(\chi_0^-V)\bigr\|_{0,p}
&\leq\|V\|_{0,p,\alpha}.
\tag{7.1.36.2}\label{eq:7-1-36-2}
\end{align}
在 $\chi_0^+V$ 的支集上，
$\alpha$ 与定义分支 $0$ 的 $\|\cdot\|_{0,p}$
所用权重一致，故第一式成立；第二式同理。
\end{proof}

现在由引理~\ref{lem:7-1-32} 估计
$\widetilde Q$ 与真正右逆的偏差。
为使分支次序清楚，记
\[
A_1:=J_{*,1}^{S}(\chi_0^-V),
\qquad A_0:=J_{*,0}^{-S}(\chi_0^+V).
\]
式~\eqref{eq:7-1-30} 给出
\[
\begin{aligned}
&\bigl\|\pi D_w\bar\partial\,
\widetilde Q_{((\mathfrak t',\ell'_T),w)}(\bar V)
-\bar V\bigr\|\\
&\quad\leq
\bigl\|\pi(J_{1,*}^{T}+J_{0,*}^{-T})
D_{(w_1,w_0)}\bar\partial Q_0(A_1\oplus A_0)
-\bar V\bigr\|
+r(T,w)\|\bar V\|\\
&\quad\leq
\bigl\|\pi\bigl(J_{1,*}^{T}A_1
+J_{0,*}^{-T}A_0\bigr)-\bar V\bigr\|
+r(T,w)\|\bar V\|,
\end{aligned}
\]
其中
$r(T,w)\leq C(e^{-cT}
+\operatorname{dist}(w,w'_0)
+\operatorname{dist}(w,w'_1))$。
第一项实际为零，因为
\[
J_{1,*}^{T}J_{*,1}^{S}(\chi_0^-V)=\chi_0^-V,
\qquad
J_{0,*}^{-T}J_{*,0}^{-S}(\chi_0^+V)=\chi_0^+V,
\]
且 $\chi_0^++\chi_0^-=1$。
这里再次用到 $\chi_0^+V$
在 $[-5T,-2]\times[0,1]$ 上为零、
$\chi_0^-V$ 在 $[2,5T]\times[0,1]$ 上为零。
\translatornote{原书把此误差简写为 $o(T)$ 并称它随 $T\to\infty$ 归零；但引理 7.1.32 还含 $w$ 与两侧基映射的距离项。对固定大小的 $w$-邻域，该收敛不能直接推出。此处保留完整误差界，先取 $T$ 足够大，再缩小 $w$ 的邻域使算子范数小于 $1$，已足以构造右逆。见疑误 A-0423。}
因此 $\pi D_w\bar\partial\,\widetilde Q$
是恒等算子的小范数扰动。
引理~\ref{lem:7-1-33} 给出
$\widetilde Q$ 的一致有界性，
再用逆算子级数即可取到满足
式~\eqref{eq:7-1-28-1}--\eqref{eq:7-1-28-2}
的 $Q_{((\mathfrak t',\ell'_T),w)}$。
\translatornote{原书此处称引理 7.1.33 已证明 $Q$ 一致有界；该引理实际估计的是尚待修正的近似右逆 $\widetilde Q$。见疑误 A-0421。}
\hfill$\square$\par

\medskip\noindent\textbf{引理~\ref{lem:7-1-29} 的证明。}
仍用该引理的情形和记号，先陈述两个辅助引理。

\begin{lemma}{7.1.37}\label{lem:7-1-37}
设 $V\in\operatorname{Ker}
D_{((\mathfrak t,\ell),w,\vec p)}\bar\partial$，
$e\in C^1_{\mathrm{ext}}(\mathfrak t)$。
照前文把 $L_e$ 识别为
$(-\infty,0]\times[0,1]$。
则存在与 $T,\delta$ 无关的常数 $c,C>0$，使
\[
|V(\tau,t)|<Ce^{-c|\tau|}\|V\|_{0,p,\alpha}.
\]
\end{lemma}

这是标准的衰减估计，亦属于引理~\ref{lem:7-1-11}
证明所用的估计。
下一个引理由伪全纯圆盘的可去奇性定理推出：
参见 \textbf{[Oh92]} 第 138 页的估计，
再将其改写成穿孔附近
$(-\infty,0]\times[0,1]$ 坐标中的形式。

\begin{lemma}{7.1.38}\label{lem:7-1-38}
设
$((\mathfrak t',\ell'_T),w,\vec p')
\in W_{k+m}^{1,p}(M,L)$
满足式~\eqref{eq:7-1-22}。\translatornote{原书引理 7.1.38 把胶合后对象的空间下标写成 $k$，并以未加撇号的 $\mathfrak t$ 标记外边；译文依本节胶合后的型别分别写 $k+m$ 和 $\mathfrak t'$。见疑误 A-0422。}
又设 $e\in C^1_{\mathrm{ext}}(\mathfrak t')$
是一条外边。将 $L_e$ 识别为
$(-\infty,0]\times[0,1]$。
则存在与 $T,\delta$ 无关的正常数
$C,C_\ell,c,c_\ell$，使
\[
\operatorname{dist}(w(\tau,t),p_i)
<Ce^{-c|\tau|},
\qquad
\|\nabla^\ell w(\tau,t)\|
<C_\ell e^{-c_\ell|\tau|}.
\]
这里 $p_i$ 是该外边对应的渐近点。
\end{lemma}

引理~\ref{lem:7-1-29}
由这两个引理及前述构造立即得到。
\hfill$\square$\par

至此，已为 $\Sigma$ 只有一个边界奇点的
$((\Sigma,\vec z),w)$ 构造仓西邻域。
边界有多个奇点时可同样处理；
内部奇点的情形可像 \textbf{[FuOn99II]} 第 3 章那样处理，
下一小节将再说明。
还可以证明这里构造的
$((\Sigma,\vec z),w)$ 的仓西邻域具有“满射性”：
它包含模空间
$\mathcal M_{k+1}(\beta)$ 中所有充分接近该点的元素。
这一满射性的证明类似
\textbf{[FuOh97]} 第 14 节或
\textbf{[FuOn99II]} 第 9 节，
也要用到引理~\ref{lem:7-1-38}。
在 \textbf{[FOOO09II]} 中，
作者还证明了一个相近但更困难情形的满射性。

结束本小节前指出，
第 7.1.2、7.1.3 小节发展出的分析结果
也可类似地用于博特--莫尔斯情形。
事实上，这里的论证立足于
\textbf{[Fuk96II]}，
该文奠定了博特--莫尔斯情形分析结果的基础。

对于某些博特--莫尔斯情形，
另有一种处理胶合分析细节的方法：
唐纳森与彼得·克朗海默（Peter Kronheimer）
在规范场论中发现了这种方法，见 \textbf{[DoKr90]}；
杜萨·麦克达夫（Dusa McDuff）与萨拉蒙
将同一方法用于闭黎曼曲面上的伪全纯曲线，
见 \textbf{[McSa94]}。
两种方法的主要区别在颈部区域的度量。
本书为颈部赋予柱面度量，
使其等距于 $[-5T,5T]\times[0,1]$
（或 $[-5T,5T]\times S^1$）；
唐纳森--克朗海默方法
（或麦克达夫--萨拉蒙方法）
则把颈部看作两个
$D^4\setminus\{0\}$
（或两个 $D^2\setminus\{0\}$）的胶合。
在 \textbf{[FuOn99II]} 中，
第一、第四署名作者采用麦克达夫--萨拉蒙方法处理球面气泡。
这里采用 \textbf{[Fuk96II]} 的方法，
因为参数化
$\mathcal M_{k+1}^{\mathrm{main}}\cong
\operatorname{Gr}_{k+1}^{+}$
与颈部区域的柱面参数化联系更直接。
